% =====================================================================
% SEBA-XAI - Introduction-only build
%
% Purpose: a clean, compilable Overleaf draft that contains ONLY the
% final professor-level Introduction and its references. All other
% sections are present as empty placeholders so the section ordering,
% numbering, and bibliography flow match the eventual full paper.
%
% This file does NOT replace main.tex. It is a sibling build target you
% can switch to in Overleaf (Menu -> Settings -> Main document ->
% main_intro_only.tex) when you want to share just the Introduction with
% your supervisor.
%
% Bibliography: references_intro_only.bib (12 keys, exactly the ones
% cited in the Introduction).
% =====================================================================

\documentclass[conference]{IEEEtran}

\usepackage{cite}
\usepackage{amsmath,amssymb}
\usepackage{graphicx}
\usepackage{booktabs}
\usepackage{url}
\usepackage{array}

\begin{document}

\title{SEBA-XAI: A Secure, Explainable, Blockchain-Audited Overlay for Inter-Agency Police Record Access Governance}

\author{
\IEEEauthorblockN{Alapati Venkat Rayudu}
\IEEEauthorblockA{
Department of Computer Science and Engineering\\
SRM University AP\\
Amaravati, India\\
Email: alapativenkatarayudu@gmail.com}
}

\maketitle

\begin{abstract}
Audit pipelines for sensitive records contain a quiet blind spot. When
the signing authority is compromised, ledger-only integrity checks,
ABAC re-execution, and Fabric-style chaincode replay all pass on a
corrupted but validly re-signed log. This paper names that scenario the
\emph{validly re-signed compromised signer} and studies it inside a
CCTNS- and ICJS-compatible secure overlay for inter-agency police
record access governance. We propose SEBA-XAI, an overlay that combines
contextual access control using RBAC, ABAC, and PBAC principles,
off-chain sensitive-record commitments, permissioned blockchain-style
audit, and logged explanation artifacts bound to audit events by hash.
The novel mechanism inside SEBA-XAI is NS-PI, an interpretable
neuro-symbolic policy-induction and drift-detection component that
operates with log-only visibility. NS-PI is evaluated side by side with
a trusted raw-attribute policy oracle baseline that assumes an
independent view of the original requests, so the paper measures the
gap between the two visibility regimes rather than presenting either as
universally best. The evaluation uses a reproducible synthetic
CCTNS/ICJS-style workload, because actual police access logs are not
publicly available. The paper does not claim live deployment, real
police records, legal compliance, or improved crime-prediction
performance.
\end{abstract}

\begin{IEEEkeywords}
Blockchain audit, access control, explainable AI, police data governance, CCTNS, ICJS, ABAC, XAI, security.
\end{IEEEkeywords}

% ---------------------------------------------------------------------
% Section I -- Introduction (the only fully written section in this build)
% ---------------------------------------------------------------------
\section{Introduction}

Audit pipelines for sensitive records share a quiet blind spot. When the
attacker controls the signing authority, every standard line of defence
collapses at once: an integrity check passes because the chain is still
internally consistent, an ABAC re-execution passes because it replays the
same recorded attributes, and a Fabric-style chaincode replay passes for
the same reason. The signer is bad, but the bytes look good. This paper
takes that scenario, names it the \emph{validly re-signed compromised
signer}, and asks a narrower question than the usual access-control paper:
under what visibility assumptions can an auditor still notice that
something is wrong, and at what cost to reviewability? The work makes that
question concrete inside a CCTNS- and ICJS-style overlay setting, but the
contribution is the joint, multi-baseline evaluation, not the choice of
domain.

The domain is what gives the question urgency. The Crime and Criminal
Tracking Network and Systems (CCTNS) runs across all 17{,}798 police
stations in India as of 1~February 2026~\cite{pib_cctns_2026}, with the
system digitising FIRs and chargesheets, hosting state-level applications,
replicating data near-real-time to a National Data Centre, and
standardising both Integrated Investigation Forms and the master codes for
states, districts, police stations, acts, and sections~\cite{pib_cctns_2026}.
Above CCTNS sits the Inter-Operable Criminal Justice System, which connects
the police pillar with courts, prisons, forensics, and prosecution through
a Data Sharing Matrix~\cite{mha_icjs_2026}. Sensitive records flow across
this fabric: FIR detail, witness and victim statements, juvenile records,
forensic reports, cybercrime complaints, case-diary text, and documents
linked to courts or prosecutors. A request from a user whose only attribute
is the broad role of ``police officer'' is not, by itself, justification
for access. The decision depends on rank, jurisdiction, case assignment,
purpose, record sensitivity, credential status, time window, emergency
context, and whether the request originates from a court, prosecutor, or
forensic workflow. Any one of these can move the same request between
legitimate and inappropriate.

\noindent\textbf{Problem statement.} Given an inter-agency request for a
sensitive police or criminal-justice record, the system should produce one
of three decisions, allow, deny, or escalate, using contextual policy
rules, and must keep an auditable and explainable record of how the
decision was reached. The record should be enough for a later, independent
reviewer to reconstruct the request, the policy version applied, the
attributes that drove the decision, the explanation produced for it, and
the approval or review event bound to it. The audit record must remain
trustworthy even when one of the system's own signing authorities is
compromised, and no single technical layer should be expected to deliver
confidentiality, correctness, and accountability at once.

The design of any such overlay rests on three layers, none of which is
novel in isolation. The first is contextual access control: NIST's guide
to Attribute-Based Access Control defines authorization in terms of
subject, object, action, and environmental attributes evaluated against
policy rules, and identifies ABAC as suitable for sharing across
organisational boundaries~\cite{nist_abac_2014}. We add policy-based
elements such as policy versioning, purpose constraints, credential
revocation, approval requirements, and superior review. The second layer
is a permissioned blockchain-style audit. Blockchain is not used here as a
privacy mechanism and is not used to store raw records. NISTIR~8403
describes blockchain-based access-control systems in terms of
decentralisation, high confidence, and tamper resistance, while also
recording the governance and scalability concerns that constrain any
honest deployment~\cite{nist_blockchain_access_2022}. Hyperledger Fabric is
the natural reference point for the permissioned, known-organisation
setting~\cite{androulaki_fabric_2018}. The overlay therefore keeps
sensitive records off-chain in agency storage and writes only minimised
metadata and commitments to the audit layer.

The third layer is where the design choice has the largest downstream
consequence. Access decisions in public-safety workflows are read by
several different people. An officer needs to know why the request went
through, was refused, or was sent for review. A supervisor needs enough
detail to act on an escalation. An auditor needs to verify that what
happened matched the declared policy. Treating explanation as a screen
shown after the decision is too weak. We treat it instead as a logged
artifact, written into the audit record alongside the decision: a decision
label, reason code, the subject/object/action/environment attributes that
mattered, the policy version, the model or rule version, counterfactual
information where it applies, and a hash that ties the explanation back to
the audit event. This stance follows Rudin's argument that high-stakes
decisions should prefer interpretable models where this is
feasible~\cite{rudin_2019}, and recent law-enforcement XAI work that
emphasises stakeholder-specific explanations, AI literacy, and the risk of
automation bias~\cite{zocholl_xai_law_enforcement_2025}.

It is worth being explicit about what this paper is not. It is not a
predictive-policing system. Ensign et al. have shown how runaway feedback
loops can arise when observed or reported crime data is itself shaped by
earlier policing activity~\cite{ensign_feedback_2018}, and earlier
criminal-justice AI work has cautioned that complex risk-prediction tools
do not always outperform simple
baselines~\cite{dressel_fairness_limits_2018}. Fairness research has also
shown that several reasonable fairness criteria become mutually
incompatible when base rates differ across
groups~\cite{chouldechova_2017,kleinberg_tradeoffs_2016}. National Crime
Records Bureau statistics, while useful as aggregate context, are not a
public, individual-level CCTNS or FIR access-decision
dataset~\cite{ncrb_2023}. We therefore make no claim to predict criminals,
identify suspects, or substitute for the judgement of an officer.

What is genuinely under-studied, given this design, is the joint
evaluation. Each pillar already has its own literature. Blockchain has
been applied to digital crime evidence and lawful chain-of-custody
workflows~\cite{kim_two_level_2021,li_lechain_2021,jeong_fabric_evidence_2020},
and auditable access-control patterns for distributed cross-organisation
processes have been studied
independently~\cite{akhtar_auditable_access_2020}. Hyperledger Fabric
combined with ABAC and encrypted off-chain storage has been studied in
generic data-sharing settings~\cite{zhao_fabric_abac_2022}, and surveys of
blockchain-based access control catalogue the design space and constrain
what counts as a novel
claim~\cite{rouhani_access_control_survey_2019,nist_blockchain_access_2022}.
More recent work on accountable and privacy-preserving blockchain access
control targets permission invisibility and access
anonymity~\cite{hu_accountable_private_ac_2024}, though it does not address
the police inter-agency workflow studied here. The XAI and fairness
literature has long flagged risks specific to criminal-justice
systems~\cite{rudin_2019,ensign_feedback_2018,zocholl_xai_law_enforcement_2025}.
None of these works, individually or together, measures what happens when
the compromised signer attacker is added to the picture. In particular,
ledger-only integrity is silent because the chain is validly signed, and
chaincode replay is silent because it operates on the corrupted canonical
output. The open question is whether anything else in a sensible overlay
notices, and what visibility assumption it must hold to do so.

The overlay we propose to answer that question is SEBA-XAI, a Secure,
Explainable, Blockchain-Audited access-governance overlay for sensitive
inter-agency police record sharing. The technical novelty of the work sits
in two places. First, we introduce NS-PI, an interpretable neuro-symbolic
policy-induction and drift-detection component that operates with
\emph{log-only visibility}: it sees nothing more than the signed decision
log and the learned policy distribution, and it raises an alarm when the
two diverge. Second, we evaluate NS-PI side by side with a trusted
raw-attribute policy oracle, which is given a stronger independent view of
the original requests and is used as a baseline for what is achievable
under that stronger assumption. The two detectors are different in trust
assumption and in mechanism, and the paper measures the gap between them
rather than presenting either as universally best. The benchmark itself is
synthetic by design, because actual police access logs and sensitive
records are not publicly available and should not be assumed in a
first-stage academic prototype.

The paper makes five initial contributions:
\begin{itemize}
  \item It introduces NS-PI, an interpretable log-only policy-drift
        detector, and the SEBA-XAI overlay in which NS-PI is one of
        several evaluated defenses.
  \item It formulates and releases a reproducible synthetic
        access-governance benchmark for CCTNS- and ICJS-style inter-agency
        record requests, with a decision space of allow, deny, or
        escalate.
  \item It introduces the validly re-signed compromised-signer attacker
        for this setting, together with an ordinary tamper catalog
        (replay, backdate, swap, collusion, revocation race, metadata
        inference).
  \item It evaluates eight defenses under explicit visibility
        assumptions, namely ledger-only integrity, ABAC- and Fabric-style
        policy re-execution~\cite{zhao_fabric_abac_2022,nist_abac_2014},
        trusted raw-attribute policy re-evaluation, and the new
        log-only NS-PI drift detector.
  \item It measures the reviewability of the overlay through audit
        reconstruction, trace completeness, decisive-attribute coverage,
        explanation stability, counterfactual coverage, and latency and
        storage overhead.
\end{itemize}

The scope of the paper is, deliberately, narrow. We do not claim live
deployment inside CCTNS or ICJS. We do not use actual police records. We
do not place raw sensitive records on-chain. We do not establish legal
compliance. We do not claim improved crime-prediction performance. The
blockchain component in the prototype is a permissioned-audit simulation,
not a live Hyperledger Fabric network. The results should be read as
reproducible benchmark evidence for a research architecture, not as
operational validation. The remainder of the paper reviews related work in
Section~II, describes the SEBA-XAI methodology and threat model in
Section~III, presents the synthetic evaluation, discusses limitations, and
outlines the future work that would be needed before any real-world pilot
could be considered.


% ---------------------------------------------------------------------
% Sections II onward: intentionally left empty in this build.
% Each \section{} preserves the eventual ordering and lets IEEEtran emit
% the expected section numbers without any prose.
% ---------------------------------------------------------------------

\section{Related Work}
% Intentionally left blank in the Introduction-only build.

\section{Threat Model}
% Intentionally left blank in the Introduction-only build.

\section{Proposed Methodology}
% Intentionally left blank in the Introduction-only build.

\section{Evaluation Setup}
% Intentionally left blank in the Introduction-only build.

\section{Results}
% Intentionally left blank in the Introduction-only build.

\section{Discussion}
% Intentionally left blank in the Introduction-only build.

\section{Limitations}
% Intentionally left blank in the Introduction-only build.

\section{Conclusion and Future Work}
% Intentionally left blank in the Introduction-only build.

% ---------------------------------------------------------------------
% References: bibliography file restricted to the 12 keys cited above.
% ---------------------------------------------------------------------
\bibliographystyle{IEEEtran}
\bibliography{references_intro_only}

\end{document}
